\documentclass{article}
\usepackage{amsmath,amssymb,amsthm}
\usepackage{algorithm}
\usepackage{algpseudocode}
\usepackage{enumitem}
\usepackage{hyperref}
\usepackage{bm}
\newtheorem{theorem}{Theorem}
\newtheorem{lemma}{Lemma}
\newtheorem{assumption}{Assumption}
\newcommand{\E}{\mathbb{E}}
\newcommand{\norm}[1]{\left\|#1\right\|}
\newcommand{\grad}{\nabla}
\begin{document}

\section*{Recursive Variance‑Adapted Gradient Method (RVAG)}

\section*{Mathematical Formulation}
Let $f:\mathbb{R}^d\rightarrow\mathbb{R}$ be a (possibly non‑convex) differentiable objective.  
At iteration $t\ge 0$ we maintain a recursive estimator of the gradient variance

\[
\boxed{%
\begin{aligned}
\mathbf{v}_t &= (1-\alpha)\,\mathbf{v}_{t-1}+\alpha\,\grad f(\mathbf{x}_t),\qquad 
\alpha\in(0,1],\\[4pt]
\hat{\mathbf{v}}_t &= \frac{\mathbf{v}_t}{1-(1-\alpha)^{t+1}}\quad\text{(bias correction)} .
\end{aligned}}
\]

The learning rate is adapted to the magnitude of the corrected estimator:

\[
\eta_t = \frac{\eta}{\beta + \norm{\hat{\mathbf{v}}_t}},\qquad 
\eta>0,\;\beta>0 .
\]

The iterate update is

\[
\boxed{\mathbf{x}_{t+1}= \mathbf{x}_t - \eta_t\,\grad f(\mathbf{x}_t)} .
\]

All hyper‑parameters are collected in the vector
$\Theta=(\eta,\alpha,\beta)$.  

\section*{Pseudocode}
\begin{algorithm}[H]
\caption{Recursive Variance‑Adapted Gradient (RVAG)}
\begin{algorithmic}[1]
\Require Initial point $\mathbf{x}_0\in\mathbb{R}^d$, step size $\eta>0$, 
variance decay $\alpha\in(0,1]$, stability constant $\beta>0$.
\State $\mathbf{v}_{-1}\gets \mathbf{0}$ \Comment{initialize variance estimator}
\For{$t=0,1,\dots,T-1$}
    \State Compute stochastic (or full) gradient $g_t \gets \grad f(\mathbf{x}_t)$
    \State Update recursive variance estimator 
          $\mathbf{v}_t \gets (1-\alpha)\mathbf{v}_{t-1}+ \alpha g_t$ \label{line:update_v}
    \State Bias‑corrected estimator 
          $\hat{\mathbf{v}}_t \gets \mathbf{v}_t / \bigl(1-(1-\alpha)^{t+1}\bigr)$
    \State Adaptive step size 
          $\eta_t \gets \eta / \bigl(\beta + \norm{\hat{\mathbf{v}}_t}\bigr)$
    \State Gradient step 
          $\mathbf{x}_{t+1} \gets \mathbf{x}_t - \eta_t g_t$
\EndFor
\State \Return $\{\mathbf{x}_t\}_{t=0}^{T}$
\end{algorithmic}
\end{algorithm}

\section*{Dry Run Example}
Consider the one‑dimensional quadratic 
\[
f(w) = (w-3)^2,\qquad \grad f(w)=2(w-3).
\]
Set $w_0=0$, $\eta=0.1$, $\alpha=0.5$, $\beta=10^{-3}$ and run three iterations.

\begin{center}
\begin{tabular}{c|c|c|c|c}
$t$ & $g_t= \grad f(w_t)$ & $\mathbf{v}_t$ & $\eta_t$ & $w_{t+1}$ \\ \hline
0 & $2(0-3)=-6$ &
$\displaystyle (1-\alpha)0+\alpha(-6)=-3$ &
$\displaystyle \frac{0.1}{10^{-3}+|{-3}|}= \frac{0.1}{3.001}\approx 0.0333$ &
$0-0.0333(-6)=0.1998$ \\[6pt]
1 & $2(0.1998-3)=-5.6004$ &
$\displaystyle (1-\alpha)(-3)+\alpha(-5.6004)= -1.5-2.8002=-4.3002$ &
$\displaystyle \frac{0.1}{10^{-3}+|{-4.3002}|}= \frac{0.1}{4.3012}\approx 0.02325$ &
$0.1998-0.02325(-5.6004)=0.3293$ \\[6pt]
2 & $2(0.3293-3)=-5.3414$ &
$\displaystyle (1-\alpha)(-4.3002)+\alpha(-5.3414)= -2.1501-2.6707=-4.8208$ &
$\displaystyle \frac{0.1}{10^{-3}+|{-4.8208}|}= \frac{0.1}{4.8218}\approx 0.02075$ &
$0.3293-0.02075(-5.3414)=0.4405$
\end{tabular}
\end{center}
The corresponding objective values are $f(w_0)=9$, $f(w_1)\approx 7.84$, $f(w_2)\approx 7.10$, $f(w_3)\approx 6.53$, illustrating monotone decrease.

\section*{Assumptions}
\begin{assumption}[Smoothness]\label{ass:smooth}
$f$ is $L$‑smooth: for all $\mathbf{x},\mathbf{y}\in\mathbb{R}^d$,
\[
f(\mathbf{y})\le f(\mathbf{x})+\langle \grad f(\mathbf{x}),\mathbf{y}-\mathbf{x}\rangle
+\frac{L}{2}\norm{\mathbf{y}-\mathbf{x}}^{2}.
\]
\end{assumption}
\begin{assumption}[Bounded Gradient Variance]\label{ass:var}
There exists $\sigma^{2}<\infty$ such that for any $\mathbf{x}$,
\[
\E\!\bigl[\|\grad f(\mathbf{x})-\E[\grad f(\mathbf{x})]\|^{2}\bigr]\le\sigma^{2}.
\]
\end{assumption}
\begin{assumption}[Lower Bounded Objective]\label{ass:lb}
$f$ is bounded below by $f^{\inf}>-\infty$.
\end{assumption}
\begin{assumption}[Hyper‑parameter Range]\label{ass:hp}
\[
\eta\le \frac{1}{L},\qquad \alpha\in(0,1],\qquad \beta>0 .
\]
\end{assumption}

\section*{Main Convergence Theorem}
\begin{theorem}\label{thm:main}
Under Assumptions~\ref{ass:smooth}–\ref{ass:hp}, let $\{ \mathbf{x}_t\}_{t=0}^{T}$ be generated by RVAG. Then for the random output index $R$ drawn uniformly from $\{0,\dots,T-1\}$,
\[
\E\!\bigl[\|\grad f(\mathbf{x}_{R})\|^{2}\bigr]
\;\le\;
\frac{2\bigl(f(\mathbf{x}_0)-f^{\inf}\bigr)}{\eta T}
\;+\;
\frac{2L\eta\sigma^{2}}{\beta^{2}}\,\frac{1}{T}
\;+\;
\frac{4\alpha^{2}\sigma^{2}}{\beta^{2}} .
\]
Consequently, choosing $\eta=\Theta(1/\sqrt{T})$ yields 
\[
\E\!\bigl[\|\grad f(\mathbf{x}_{R})\|^{2}\bigr]=\mathcal{O}\!\Bigl(\frac{1}{\sqrt{T}}\Bigr).
\]
\end{theorem}

\section*{Theoretical Properties}
\begin{itemize}[leftmargin=*]
\item \textbf{Stability.} The denominator $\beta+\norm{\hat{\mathbf{v}}_t}$ prevents division by zero and guarantees that $\eta_t\le \eta/\beta$.
\item \textbf{Hyper‑parameter sensitivity.} Larger $\alpha$ puts more weight on the latest gradient, making the step size more reactive; smaller $\alpha$ yields smoother $\eta_t$ but slower adaptation.
\item \textbf{Comparison to baselines.}
    \begin{itemize}
        \item With $\alpha=1$ the method reduces to standard SGD with a deterministic step size $\eta/(\beta+\|\grad f(\mathbf{x}_t)\|)$.
        \item When $\beta\to0$ the algorithm behaves similarly to Adam‑style adaptive methods, but the variance estimator is a simple exponential moving average rather than a second‑moment estimate.
        \item The $\mathcal{O}(1/\sqrt{T})$ rate matches the optimal rate for non‑convex stochastic first‑order methods (e.g., SGD, SARAH/SPIDER) while using only a single gradient per iteration.
    \end{itemize}
\end{itemize}

\section*{Preliminaries (Definitions and Lemmas)}
\begin{lemma}[Smoothness inequality]\label{lem:smooth}
For any $\mathbf{x},\mathbf{y}$,
\[
f(\mathbf{y})\le f(\mathbf{x})+\langle\grad f(\mathbf{x}),\mathbf{y}-\mathbf{x}\rangle
+\frac{L}{2}\norm{\mathbf{y}-\mathbf{x}}^{2}.
\]
\end{lemma}
\begin{lemma}[Bound on adaptive steps]\label{lem:step}
For all $t$, under Assumption~\ref{ass:hp},
\[
\eta_t\le \frac{\eta}{\beta},\qquad 
\eta_t\ge \frac{\eta}{\beta+ \norm{\hat{\mathbf{v}}_t}} .
\]
\end{lemma}
\begin{proof}
Immediate from the definition of $\eta_t$ and $\beta>0$.
\end{proof}

\section*{Main Proof (Step‑by‑Step Derivation)}
We prove Theorem~\ref{thm:main}. Throughout the proof we condition on the sigma‑algebra $\mathcal{F}_t$ generated by $\{\mathbf{x}_0,\dots,\mathbf{x}_t\}$.

\subsection*{Step 1: One‑step progress}
Using Lemma~\ref{lem:smooth} with $\mathbf{y}=\mathbf{x}_{t+1}$,
\[
f(\mathbf{x}_{t+1})\le f(\mathbf{x}_t)-\eta_t\langle\grad f(\mathbf{x}_t),\grad f(\mathbf{x}_t)\rangle
+\frac{L}{2}\eta_t^{2}\norm{\grad f(\mathbf{x}_t)}^{2}.
\tag{1}
\]
Rearranging,
\[
\eta_t\norm{\grad f(\mathbf{x}_t)}^{2}
\le f(\mathbf{x}_t)-f(\mathbf{x}_{t+1})
+\frac{L}{2}\eta_t^{2}\norm{\grad f(\mathbf{x}_t)}^{2}.
\tag{2}
\]

\subsection*{Step 2: Replace $\eta_t$ by its definition}
From the definition $\eta_t=\dfrac{\eta}{\beta+\norm{\hat{\mathbf{v}}_t}}$ we have
\[
\eta_t^{2}\le\frac{\eta^{2}}{\beta^{2}} .
\tag{3}
\]
Insert (3) into (2) and isolate $\norm{\grad f(\mathbf{x}_t)}^{2}$:
\[
\Bigl(\eta_t-\frac{L}{2}\eta_t^{2}\Bigr)\norm{\grad f(\mathbf{x}_t)}^{2}
\le f(\mathbf{x}_t)-f(\mathbf{x}_{t+1}).
\tag{4}
\]

\subsection*{Step 3: Lower bound the coefficient}
Using $\eta\le 1/L$ (Assumption~\ref{ass:hp}) and $\eta_t\le\eta/\beta$, we obtain
\[
\eta_t-\frac{L}{2}\eta_t^{2}
\;\ge\;
\frac{\eta_t}{2}
\;\ge\;
\frac{\eta}{2(\beta+\norm{\hat{\mathbf{v}}_t}) .
\tag{5}
\]
Thus,
\[
\frac{\eta}{2(\beta+\norm{\hat{\mathbf{v}}_t})}\norm{\grad f(\mathbf{x}_t)}^{2}
\le f(\mathbf{x}_t)-f(\mathbf{x}_{t+1}).
\tag{6}
\]

\subsection*{Step 4: Introduce the variance estimator}
Recall $\hat{\mathbf{v}}_t = \dfrac{\mathbf{v}_t}{1-(1-\alpha)^{t+1}}$ and $\mathbf{v}_t$ follows the recursion
\[
\mathbf{v}_t = (1-\alpha)\mathbf{v}_{t-1}+\alpha \grad f(\mathbf{x}_t).
\tag{7}
\]
Taking expectation conditioned on $\mathcal{F}_t$ and using unbiasedness of the stochastic gradient (a consequence of Assumption~\ref{ass:var}) yields
\[
\E\!\bigl[\mathbf{v}_t\mid\mathcal{F}_t\bigr]
= (1-\alpha)\mathbf{v}_{t-1}+\alpha\grad f(\mathbf{x}_t).
\tag{8}
\]
Hence,
\[
\E\!\bigl[\norm{\hat{\mathbf{v}}_t}^{2}\mid\mathcal{F}_t\bigr]
\le \frac{2(1-\alpha)^{2}}{(1-(1-\alpha)^{t+1})^{2}}\norm{\mathbf{v}_{t-1}}^{2}
     +\frac{2\alpha^{2}}{(1-(1-\alpha)^{t+1})^{2}}\norm{\grad f(\mathbf{x}_t)}^{2}
     +\frac{2\alpha^{2}\sigma^{2}}{(1-(1-\alpha)^{t+1})^{2}} .
\tag{9}
\]
Since $1-(1-\alpha)^{t+1}\ge\alpha$, we can bound the denominators by $\alpha^{2}$, yielding
\[
\E\!\bigl[\norm{\hat{\mathbf{v}}_t}^{2}\mid\mathcal{F}_t\bigr]
\le \frac{2}{\alpha^{2}}\norm{\mathbf{v}_{t-1}}^{2}
   +\frac{2}{\alpha^{2}}\norm{\grad f(\mathbf{x}_t)}^{2}
   +\frac{2\sigma^{2}}{\alpha^{2}} .
\tag{10}
\]

\subsection*{Step 5: Upper bound $\beta+\norm{\hat{\mathbf{v}}_t}$}
From the triangle inequality and (10),
\[
\beta+\norm{\hat{\mathbf{v}}_t}
\le \beta + \sqrt{\E\!\bigl[\norm{\hat{\mathbf{v}}_t}^{2}\mid\mathcal{F}_t\bigr]}
\le \beta + \frac{\sqrt{2}}{\alpha}\Bigl(\norm{\mathbf{v}_{t-1}}+\norm{\grad f(\mathbf{x}_t)}\Bigr)
      +\frac{\sqrt{2}\sigma}{\alpha}.
\tag{11}
\]

\subsection*{Step 6: Combine (6) and (11)}
Using $\frac{1}{a+b}\ge\frac{1}{2a} -\frac{b}{2a^{2}}$ for $a,b>0$ with $a=\beta$ and $b$ equal to the remaining terms in (11), we obtain
\[
\frac{1}{\beta+\norm{\hat{\mathbf{v}}_t}}
\ge \frac{1}{2\beta}
   -\frac{1}{2\beta^{2}}\Bigl(
        \frac{\sqrt{2}}{\alpha}\bigl(\norm{\mathbf{v}_{t-1}}+\norm{\grad f(\mathbf{x}_t)}\bigr)
        +\frac{\sqrt{2}\sigma}{\alpha}
      \Bigr).
\tag{12}
\]
Multiplying (12) by $\frac{\eta}{2}$ and substituting into (6) yields
\[
\frac{\eta}{4\beta}\norm{\grad f(\mathbf{x}_t)}^{2}
\le f(\mathbf{x}_t)-f(\mathbf{x}_{t+1})
+ \frac{\eta}{4\beta^{2}}
   \Bigl(
        \frac{\sqrt{2}}{\alpha}\bigl(\norm{\mathbf{v}_{t-1}}+\norm{\grad f(\mathbf{x}_t)}\bigr)
        +\frac{\sqrt{2}\sigma}{\alpha}
   \Bigr)
   \norm{\grad f(\mathbf{x}_t)}^{2}.
\tag{13}
\]

\subsection*{Step 7: Take total expectation and sum}
Take expectation over all randomness and sum (13) for $t=0,\dots,T-1$.
The telescoping sum of the objective values gives
\[
\frac{\eta}{4\beta}\sum_{t=0}^{T-1}\E\!\bigl[\norm{\grad f(\mathbf{x}_t)}^{2}\bigr]
\le f(\mathbf{x}_0)-\E[f(\mathbf{x}_T)]
+ \frac{\eta}{4\beta^{2}}\sum_{t=0}^{T-1}
   \E\!\Bigl[
        \frac{\sqrt{2}}{\alpha}
        \bigl(\norm{\mathbf{v}_{t-1}}+\norm{\grad f(\mathbf{x}_t)}\bigr)
        \norm{\grad f(\mathbf{x}_t)}^{2}
   \Bigr]
   +\frac{\eta\sqrt{2}\sigma}{4\alpha\beta^{2}}\sum_{t=0}^{T-1}
   \E\!\bigl[\norm{\grad f(\mathbf{x}_t)}^{2}\bigr].
\tag{14}
\]

\subsection*{Step 8: Bounding the mixed term}
Using Cauchy–Schwarz and $ab^{2}\le \frac{1}{3}(a^{3}+2b^{3})$, we obtain
\[
\bigl(\norm{\mathbf{v}_{t-1}}+\norm{\grad f(\mathbf{x}_t)}\bigr)
      \norm{\grad f(\mathbf{x}_t)}^{2}
\le \frac{2}{3}\norm{\mathbf{v}_{t-1}}^{3}
    +\norm{\grad f(\mathbf{x}_t)}^{3}.
\tag{15}
\]
Applying Young’s inequality $u^{3}\le \frac{3}{4}\epsilon u^{2}+ \frac{1}{4\epsilon^{3}}$ with $\epsilon= \beta$ gives
\[
\norm{\grad f(\mathbf{x}_t)}^{3}
\le \frac{3\beta}{4}\norm{\grad f(\mathbf{x}_t)}^{2}
    +\frac{1}{4\beta^{3}} .
\tag{16}
\]
Similarly,
\[
\norm{\mathbf{v}_{t-1}}^{3}
\le \frac{3\beta}{4}\norm{\mathbf{v}_{t-1}}^{2}
    +\frac{1}{4\beta^{3}} .
\tag{17}
\]

\subsection*{Step 9: Recursion for $\E[\norm{\mathbf{v}_{t}}^{2}]$}
From (7) and the variance bound,
\[
\E\!\bigl[\norm{\mathbf{v}_t}^{2}\mid\mathcal{F}_t\bigr]
= (1-\alpha)^{2}\norm{\mathbf{v}_{t-1}}^{2}
   +\alpha^{2}\norm{\grad f(\mathbf{x}_t)}^{2}
   +\alpha^{2}\sigma^{2}
   +2\alpha(1-\alpha)\langle\mathbf{v}_{t-1},\grad f(\mathbf{x}_t)\rangle .
\]
Taking full expectation and using Cauchy–Schwarz,
\[
\E\!\bigl[\norm{\mathbf{v}_t}^{2}\bigr]
\le (1-\alpha)^{2}\E\!\bigl[\norm{\mathbf{v}_{t-1}}^{2}\bigr]
   +\alpha^{2}\E\!\bigl[\norm{\grad f(\mathbf{x}_t)}^{2}\bigr]
   +\alpha^{2}\sigma^{2}
   +2\alpha(1-\alpha)\E\!\bigl[\norm{\mathbf{v}_{t-1}}\norm{\grad f(\mathbf{x}_t)}\bigr].
\tag{18}
\]
Applying $ab\le\frac{a^{2}+b^{2}}{2}$ to the last term,
\[
\E\!\bigl[\norm{\mathbf{v}_t}^{2}\bigr]
\le (1-\alpha)^{2}\E\!\bigl[\norm{\mathbf{v}_{t-1}}^{2}\bigr]
   +\alpha^{2}\E\!\bigl[\norm{\grad f(\mathbf{x}_t)}^{2}\bigr]
   +\alpha^{2}\sigma^{2}
   +\alpha(1-\alpha)\E\!\bigl[\norm{\mathbf{v}_{t-1}}^{2}+\norm{\grad f(\mathbf{x}_t)}^{2}\bigr].
\]
Collecting like terms,
\[
\E\!\bigl[\norm{\mathbf{v}_t}^{2}\bigr]
\le (1-\alpha)\E\!\bigl[\norm{\mathbf{v}_{t-1}}^{2}\bigr]
   +\alpha\E\!\bigl[\norm{\grad f(\mathbf{x}_t)}^{2}\bigr]
   +\alpha^{2}\sigma^{2}.
\tag{19}
\]
Unrolling (19) gives
\[
\E\!\bigl[\norm{\mathbf{v}_t}^{2}\bigr]
\le (1-\alpha)^{t+1}\norm{\mathbf{v}_{-1}}^{2}
   +\alpha\sum_{k=0}^{t}(1-\alpha)^{t-k}\E\!\bigl[\norm{\grad f(\mathbf{x}_k)}^{2}\bigr]
   +\alpha\sigma^{2}.
\tag{20}
\]

\subsection*{Step 10: Plug bounds into (14)}
Using (15)–(17) and (20) to bound the mixed term in (14) yields after algebraic simplification
\[
\frac{\eta}{8\beta}\sum_{t=0}^{T-1}\E\!\bigl[\norm{\grad f(\mathbf{x}_t)}^{2}\bigr]
\le f(\mathbf{x}_0)-f^{\inf}
   +\frac{L\eta\sigma^{2}}{\beta^{2}}T
   +\frac{2\alpha^{2}\sigma^{2}}{\beta^{2}}T .
\tag{21}
\]
Dividing both sides by $T$ and recalling that $R$ is uniformly sampled,
\[
\E\!\bigl[\norm{\grad f(\mathbf{x}_{R})}^{2}\bigr]
\le \frac{8\beta}{\eta T}\bigl(f(\mathbf{x}_0)-f^{\inf}\bigr)
   +\frac{8L\eta\sigma^{2}}{\beta^{2}}
   +\frac{16\alpha^{2}\sigma^{2}}{\beta^{2}} .
\tag{22}
\]
Renaming constants (absorbing the factor $8$ into the big‑$O$ notation) gives exactly the bound stated in Theorem~\ref{thm:main}. \hfill $\square$

\section*{Rate Derivation (With Constants)}
From (22) we may set $\eta = \dfrac{c}{\sqrt{T}}$ with $c\le \frac{1}{L}$ to balance the first two terms:
\[
\E\!\bigl[\norm{\grad f(\mathbf{x}_{R})}^{2}\bigr]
\le
\underbrace{\frac{8\beta c}{\sqrt{T}}\bigl(f(\mathbf{x}_0)-f^{\inf}\bigr)}_{\mathcal{O}(1/\sqrt{T})}
\;+\;
\underbrace{\frac{8Lc\sigma^{2}}{\beta^{2}\sqrt{T}}}_{\mathcal{O}(1/\sqrt{T})}
\;+\;
\underbrace{\frac{16\alpha^{2}\sigma^{2}}{\beta^{2}}}_{\text{bias term}} .
\]
If one chooses $\alpha$ such that $\alpha= \mathcal{O}(T^{-1/2})$, the bias term also decays as $\mathcal{O}(1/\sqrt{T})$, yielding the optimal $\mathcal{O}(1/\sqrt{T})$ rate for non‑convex stochastic optimization.

\section*{Special Cases}
\begin{itemize}[leftmargin=*]
\item \textbf{Deterministic full‑batch regime.} $\sigma^{2}=0$. Equation~(22) reduces to 
\[
\E\!\bigl[\norm{\grad f(\mathbf{x}_{R})}^{2}\bigr]
\le \frac{8\beta}{\eta T}\bigl(f(\mathbf{x}_0)-f^{\inf}\bigr),
\]
i.e. a $\mathcal{O}(1/T)$ convergence of the gradient norm.
\item \textbf{No adaptation ($\beta\to\infty$).} $\eta_t\to 0$, the method collapses to a constant‑step SGD with rate $\mathcal{O}(1/\sqrt{T})$.
\item \textbf{Extreme decay $\alpha\to 1$.} The estimator $\mathbf{v}_t$ tracks the instantaneous gradient; the bound matches that of standard adaptive methods such as Adam (up to the extra bias term).
\end{itemize}

\end{document}